% Source: https://artofproblemsolving.com/community/c3h2808194_cool_identity_of_a_parabola
% Stable identity: aops-topic-2808194-q01
The slope of $PQ$ is $\dfrac{x_1^2-x_2^2}{x_1+x_2}=x_1-x_2$. The line through $P$ is $y=x_1^2+(x_1-x_2)(x-x_1)$; at $x=0$ this gives $y=x_1^2-x_1^2+x_1x_2=x_1x_2$.
\begin{asy}
usepackage("amsmath");
size(5cm);
defaultpen(fontsize(10pt));
import graph;real f(real x){return x^2;}draw(graph(f,-2.3,2.6));draw((-2.6,0)--(2.9,0),Arrow);draw((0,-0.5)--(0,6.8),Arrow);
pair P=(2.5,6.25),Q=(-1.6,2.56);draw(Q--P);dot(P);dot(Q);dot((0,4));label("$P(x_1,x_1^2)$",P,E);label("$Q(-x_2,x_2^2)$",Q,W);label("$(0,x_1x_2)$",(0,4),SE);
\end{asy}

Since $P$ and $Q$ are on opposite sides of the $y$-axis, the segment does cross it, at $(0,x_1x_2)$. $\blacksquare$
